\documentclass[11pt,a4paper]{article}

\usepackage{amsmath,amssymb,amsthm}
\usepackage{graphicx}
\usepackage{hyperref}
\usepackage{algorithm}
\usepackage{algpseudocode}
\usepackage{booktabs}
\usepackage{geometry}
\usepackage{authblk}

\geometry{margin=1in}

\theoremstyle{definition}
\newtheorem{definition}{Definition}
\newtheorem{theorem}{Theorem}
\newtheorem{lemma}{Lemma}
\newtheorem{proposition}{Proposition}
\newtheorem{corollary}{Corollary}

\title{Yield Aggregation Strategies: \\
Auto-Compounding and Strategy Vaults on Lux}

\author{Zach Kelling}
\affil{Hanzo Industries \\ Lux Industries \\ Zoo Labs Foundation}
\date{2023}

\begin{document}

\maketitle

\begin{abstract}
We present a mathematical framework for yield aggregation protocols on the Lux blockchain. Yield aggregators optimize returns by automating compound interest harvesting, rebalancing across yield sources, and implementing sophisticated farming strategies. We derive optimal compounding frequencies, analyze gas-adjusted returns, and formalize strategy vault architectures. Our framework demonstrates that automated yield optimization significantly outperforms manual management, particularly in high-gas environments where batched operations provide economies of scale. We provide implementation guidelines and risk analysis for production deployment.
\end{abstract}

\section{Introduction}

Decentralized finance offers numerous yield-generating opportunities: lending protocols, liquidity provision, staking, and farming incentives. However, maximizing returns requires:

\begin{enumerate}
\item Frequent compounding of earned rewards
\item Rebalancing across multiple yield sources
\item Gas-efficient execution of complex strategies
\end{enumerate}

Yield aggregators automate these operations, pooling user deposits to achieve economies of scale and implementing optimized strategies programmatically.

\section{Compound Interest Theory}

\subsection{Continuous vs. Discrete Compounding}

\begin{definition}[Annual Percentage Rate]
The nominal APR $r$ represents interest earned without compounding:
\begin{equation}
\text{Value}(t) = P \cdot (1 + r \cdot t)
\end{equation}
\end{definition}

\begin{definition}[Annual Percentage Yield]
The APY accounts for compounding $n$ times per year:
\begin{equation}
\text{APY} = \left(1 + \frac{r}{n}\right)^n - 1
\end{equation}
\end{definition}

\begin{theorem}[Continuous Compounding Limit]
As compounding frequency approaches infinity:
\begin{equation}
\lim_{n \to \infty} \left(1 + \frac{r}{n}\right)^n = e^r
\end{equation}
Thus, continuous compounding yields:
\begin{equation}
\text{APY}_{\text{continuous}} = e^r - 1
\end{equation}
\end{theorem}

\begin{proof}
By definition of the exponential function:
\begin{equation}
e^x = \lim_{n \to \infty} \left(1 + \frac{x}{n}\right)^n
\end{equation}
\end{proof}

\subsection{Compounding Benefit}

\begin{proposition}[Compounding Gain]
The benefit of compounding $n$ times vs. no compounding:
\begin{equation}
\text{Gain}(n, r) = \left(1 + \frac{r}{n}\right)^n - (1 + r)
\end{equation}
\end{proposition}

\begin{table}[h]
\centering
\begin{tabular}{|c|c|c|}
\hline
Compounding & APR = 10\% & APR = 100\% \\
\hline
None (1x/year) & 10.00\% & 100.00\% \\
Monthly (12x) & 10.47\% & 161.30\% \\
Daily (365x) & 10.52\% & 171.46\% \\
Continuous & 10.52\% & 171.83\% \\
\hline
\end{tabular}
\caption{APY by compounding frequency}
\end{table}

\section{Optimal Compounding Frequency}

\subsection{Gas-Adjusted Returns}

\begin{definition}[Net APY]
For compounding cost $G$ (in deposit currency) per operation:
\begin{equation}
\text{Net APY}(n) = \left(1 + \frac{r}{n}\right)^n - 1 - \frac{n \cdot G}{P}
\end{equation}
where $P$ is principal.
\end{definition}

\begin{theorem}[Optimal Compounding Frequency]
The optimal compounding frequency $n^*$ maximizes net APY:
\begin{equation}
n^* = \underset{n \in \mathbb{Z}^+}{\arg\max} \left[\left(1 + \frac{r}{n}\right)^n - \frac{n \cdot G}{P}\right]
\end{equation}
\end{theorem}

\begin{proposition}[Optimal Frequency Approximation]
For continuous approximation:
\begin{equation}
n^* \approx \sqrt{\frac{r \cdot P}{G}}
\end{equation}
\end{proposition}

\begin{proof}
Maximize $f(n) = e^r - n \cdot G/P$ with respect to continuous compounding benefit approaching $e^r$. The marginal benefit of compounding decreases as $n$ increases, while marginal cost is constant at $G/P$. Setting marginal benefit equal to marginal cost:
\begin{equation}
\frac{d}{dn}\left[\left(1 + \frac{r}{n}\right)^n\right] \approx \frac{G}{P}
\end{equation}
For small $r/n$, this yields $n^* \approx \sqrt{rP/G}$.
\end{proof}

\subsection{Scale Economies}

\begin{corollary}[Pooled Compounding Benefit]
For a pool of $N$ depositors with total principal $P_{\text{total}}$:
\begin{equation}
n^*_{\text{pool}} = \sqrt{\frac{r \cdot P_{\text{total}}}{G}} = \sqrt{N} \cdot n^*_{\text{individual}}
\end{equation}
Pooling enables $\sqrt{N}$ times more frequent compounding.
\end{corollary}

\section{Strategy Vault Architecture}

\subsection{Vault Definition}

\begin{definition}[Strategy Vault]
A vault $V$ is defined by:
\begin{itemize}
\item Deposit asset $A$
\item Strategy contract $S$
\item Share token $T$
\item Performance fee $\phi$
\item Management fee $\mu$
\end{itemize}
\end{definition}

\subsection{Share Accounting}

\begin{definition}[Price Per Share]
The value of each vault share:
\begin{equation}
\text{PPS} = \frac{\text{Total Assets}}{\text{Total Shares}}
\end{equation}
\end{definition}

\begin{algorithm}
\caption{Deposit into Vault}
\begin{algorithmic}[1]
\Require User $u$, deposit amount $A$
\Ensure Shares minted to user
\State $\text{assetsBefore} \gets \text{vault.totalAssets()}$
\State $\text{sharesBefore} \gets \text{vault.totalSupply()}$
\State \textbf{transfer}(asset, $u$, vault, $A$)
\If{$\text{sharesBefore} = 0$}
\State $\text{shares} \gets A$ \Comment{Initial deposit}
\Else
\State $\text{shares} \gets A \cdot \text{sharesBefore} / \text{assetsBefore}$
\EndIf
\State \textbf{mint}(shareToken, $u$, shares)
\State \Return shares
\end{algorithmic}
\end{algorithm}

\begin{algorithm}
\caption{Withdraw from Vault}
\begin{algorithmic}[1]
\Require User $u$, share amount $S$
\Ensure Assets returned to user
\State $\text{totalAssets} \gets \text{vault.totalAssets()}$
\State $\text{totalShares} \gets \text{vault.totalSupply()}$
\State $\text{assets} \gets S \cdot \text{totalAssets} / \text{totalShares}$
\State \textbf{burn}(shareToken, $u$, $S$)
\State \textbf{transfer}(asset, vault, $u$, assets)
\State \Return assets
\end{algorithmic}
\end{algorithm}

\subsection{Fee Structure}

\begin{definition}[Performance Fee]
Fee on profits since last harvest:
\begin{equation}
\text{Fee}_{\text{perf}} = \phi \cdot \max(\text{TotalAssets}_t - \text{TotalAssets}_{t-1}, 0)
\end{equation}
Typical: $\phi = 10\%$ to $20\%$.
\end{definition}

\begin{definition}[Management Fee]
Annual fee on assets under management:
\begin{equation}
\text{Fee}_{\text{mgmt}} = \mu \cdot \text{TotalAssets} \cdot \Delta t
\end{equation}
Typical: $\mu = 2\%$ annually.
\end{definition}

\section{Yield Strategies}

\subsection{Single-Asset Strategies}

\begin{definition}[Lending Strategy]
Deposit asset into lending protocol, earn interest:
\begin{equation}
\text{APY} = r_{\text{supply}} \cdot U \cdot (1 - \rho)
\end{equation}
where $U$ is utilization and $\rho$ is reserve factor.
\end{definition}

\begin{definition}[Staking Strategy]
Stake asset in proof-of-stake or governance:
\begin{equation}
\text{APY} = \frac{\text{Emissions}_{\text{annual}}}{\text{Total Staked}}
\end{equation}
\end{definition}

\subsection{LP Strategies}

\begin{definition}[LP Farm Strategy]
\begin{enumerate}
\item Deposit paired assets into AMM
\item Stake LP tokens in farm
\item Harvest reward tokens
\item Swap rewards to base assets
\item Compound into more LP
\end{enumerate}
\end{definition}

\begin{proposition}[LP Strategy Returns]
Total APY for LP farming:
\begin{equation}
\text{APY}_{\text{total}} = \text{APY}_{\text{fees}} + \text{APY}_{\text{rewards}} - \text{IL}
\end{equation}
where IL is impermanent loss.
\end{proposition}

\subsection{Leveraged Strategies}

\begin{definition}[Recursive Borrowing]
Leverage yield by recursive deposit-borrow cycles:
\begin{enumerate}
\item Deposit collateral
\item Borrow against collateral
\item Deposit borrowed funds as more collateral
\item Repeat until leverage target reached
\end{enumerate}
\end{definition}

\begin{theorem}[Leverage Multiplier]
For loan-to-value $\lambda$ and $n$ cycles:
\begin{equation}
\text{Effective Exposure} = \sum_{i=0}^{n} \lambda^i = \frac{1 - \lambda^{n+1}}{1 - \lambda}
\end{equation}
As $n \to \infty$:
\begin{equation}
\text{Max Leverage} = \frac{1}{1 - \lambda}
\end{equation}
\end{theorem}

\begin{corollary}[Leveraged Yield]
Net APY with leverage:
\begin{equation}
\text{APY}_{\text{levered}} = L \cdot r_{\text{supply}} - (L - 1) \cdot r_{\text{borrow}}
\end{equation}
where $L$ is leverage multiplier.
\end{corollary}

Profitable when $r_{\text{supply}} > r_{\text{borrow}} \cdot \frac{L-1}{L}$.

\section{Strategy Selection}

\subsection{Risk-Adjusted Returns}

\begin{definition}[Sharpe Ratio]
Risk-adjusted return metric:
\begin{equation}
\text{Sharpe} = \frac{\mathbb{E}[R] - r_f}{\sigma_R}
\end{equation}
where $r_f$ is risk-free rate and $\sigma_R$ is return volatility.
\end{definition}

\begin{definition}[Strategy Risk Score]
Composite risk assessment:
\begin{equation}
\text{Risk} = w_1 \cdot R_{\text{smart contract}} + w_2 \cdot R_{\text{oracle}} + w_3 \cdot R_{\text{liquidity}} + w_4 \cdot R_{\text{IL}}
\end{equation}
\end{definition}

\subsection{Strategy Allocation}

\begin{definition}[Mean-Variance Optimization]
For $n$ strategies with expected returns $\mu$ and covariance $\Sigma$:
\begin{equation}
\max_w \quad w^T \mu - \frac{\gamma}{2} w^T \Sigma w
\end{equation}
subject to $\sum_i w_i = 1$, $w_i \geq 0$.
\end{definition}

\section{Auto-Compounding Implementation}

\subsection{Harvest Mechanism}

\begin{algorithm}
\caption{Harvest and Compound}
\begin{algorithmic}[1]
\Require Vault $V$, strategy $S$
\Ensure Rewards compounded
\State $\text{rewards} \gets S.\text{claimRewards}()$
\For{each reward token $r$}
\State $\text{baseAmount} \gets \text{swap}(r, \text{baseAsset})$
\EndFor
\State $\text{profit} \gets \text{baseAmount}$
\State $\text{fee} \gets \text{profit} \cdot \phi$
\State \textbf{mint}(shares, treasury, fee / PPS)
\State $S.\text{deposit}(\text{profit} - \text{fee})$
\State \textbf{emit} Harvested(profit, fee)
\end{algorithmic}
\end{algorithm}

\subsection{Keeper Incentives}

\begin{definition}[Keeper Reward]
Keepers who trigger harvest receive:
\begin{equation}
\text{Reward} = \min(\text{Gas} \cdot (1 + \text{premium}), \text{maxReward})
\end{equation}
\end{definition}

\begin{proposition}[Harvest Profitability]
Harvest is profitable for keepers when:
\begin{equation}
\text{Pending Rewards} > \frac{\text{Gas Cost}}{1 + \text{premium}}
\end{equation}
\end{proposition}

\section{Risk Management}

\subsection{Strategy Limits}

\begin{definition}[Debt Ratio]
Fraction of vault assets allocated to strategy:
\begin{equation}
\text{Debt Ratio} = \frac{\text{Strategy Debt}}{\text{Total Assets}}
\end{equation}
\end{definition}

\begin{definition}[Strategy Parameters]
\begin{itemize}
\item $\text{debtRatio}_{\max}$: Maximum allocation (e.g., 90\%)
\item $\text{minDebtPerHarvest}$: Minimum profit to trigger harvest
\item $\text{maxDebtPerHarvest}$: Maximum single harvest size
\end{itemize}
\end{definition}

\subsection{Emergency Procedures}

\begin{algorithm}
\caption{Emergency Withdrawal}
\begin{algorithmic}[1]
\Require Guardian authorization
\Ensure All funds returned to vault
\State $S.\text{emergencyExit}()$
\State $S.\text{setDebtRatio}(0)$
\State \textbf{emit} EmergencyExit($S$)
\end{algorithmic}
\end{algorithm}

\subsection{Slippage Protection}

\begin{definition}[Slippage Bounds]
For swaps during harvest:
\begin{equation}
\text{minAmountOut} = \text{expectedAmount} \cdot (1 - \text{slippageTolerance})
\end{equation}
Typical tolerance: 0.5\% to 2\%.
\end{definition}

\section{Implementation on Lux}

\subsection{Architecture Overview}

\begin{enumerate}
\item \textbf{Vault Contract}: Manages deposits, withdrawals, share accounting
\item \textbf{Strategy Contracts}: Implement yield-generating logic
\item \textbf{Registry}: Tracks approved vaults and strategies
\item \textbf{Keeper Network}: Automates harvest triggers
\item \textbf{Oracle Integration}: Price feeds for swap routing
\end{enumerate}

\subsection{Gas Optimization}

\begin{itemize}
\item \textbf{Batched Harvests}: Multiple strategies in single transaction
\item \textbf{Lazy Accounting}: Update balances only on user interaction
\item \textbf{Minimal Storage}: Compute values on-the-fly where possible
\end{itemize}

\subsection{Multi-Chain Considerations}

Lux's subnet architecture enables:
\begin{itemize}
\item Strategy deployment on optimized subnets
\item Cross-subnet yield aggregation
\item Subnet-specific gas optimization
\end{itemize}

\section{Performance Analysis}

\subsection{Compounding Impact}

\begin{theorem}[Aggregator Advantage]
For base APR $r$, gas cost $G$, and pool size $P$:
\begin{equation}
\text{APY}_{\text{aggregator}} - \text{APY}_{\text{manual}} \approx \frac{r^2}{2} - \frac{G}{P}\sqrt{\frac{r \cdot P}{G}}
\end{equation}
\end{theorem}

For typical parameters ($r = 50\%$, $G = \$10$, $P = \$1M$):
\begin{equation}
\text{Advantage} \approx 12.5\% - 0.07\% = 12.43\%
\end{equation}

\section{Conclusion}

Yield aggregation provides significant advantages over manual yield farming:

\begin{enumerate}
\item \textbf{Optimal Compounding}: Mathematically-derived harvest frequency
\item \textbf{Gas Efficiency}: Pooled operations reduce per-user costs
\item \textbf{Strategy Diversification}: Risk-adjusted allocation across sources
\item \textbf{Automation}: Continuous optimization without user intervention
\end{enumerate}

Lux's high throughput and low fees make it ideal for yield aggregation, enabling frequent compounding with minimal cost overhead.

\begin{thebibliography}{9}
\bibitem{yearn}
Yearn Finance. (2020). yVaults: Automated Yield Generation. Yearn Documentation.

\bibitem{convex}
Convex Finance. (2021). Boosted Curve Staking. Convex Whitepaper.

\bibitem{beefy}
Beefy Finance. (2020). Multi-chain Yield Optimizer. Beefy Documentation.

\bibitem{harvest}
Harvest Finance. (2020). Yield Farming Made Easy. Harvest Whitepaper.

\bibitem{pickle}
Pickle Finance. (2020). Yield Aggregation Protocol. Pickle Documentation.
\end{thebibliography}

\end{document}
